\unitchapter{Paper 1 · Unit 2}{Research Aptitude}{p1-02}

\textbf{Expected questions: \textasciitilde5 (10 marks) · Target: 5/5}

\begin{sylbox}

\begin{itemize}
\tightlist
\item[$\square$]
  Research: meaning, types, characteristics; positivism \& post-positivistic approach
\item[$\square$]
  Methods: experimental, descriptive, historical, qualitative, quantitative
\item[$\square$]
  Steps of research
\item[$\square$]
  Thesis \& article writing: format and styles of referencing
\item[$\square$]
  Application of ICT in research
\item[$\square$]
  Research ethics
\end{itemize}

\end{sylbox}

\needdiagram{figures/p1-02-00}{130pt}

\hypertarget{p1-02--1-meaning--characteristics}{%
\section{1. Meaning \& Characteristics}\label{p1-02--1-meaning--characteristics}}

Research = \textbf{systematic, objective, controlled, empirical, critical} investigation to discover new facts or verify old ones.
(\emph{Re + search} = search again.) Characteristics: systematic, logical, empirical, replicable, reductive, controlled, generalisable.

\needspace{100pt}

\hypertarget{p1-02--2-types-of-research}{%
\section{2. Types of Research}\label{p1-02--2-types-of-research}}

\diagramhere{figures/p1-02-00}

\begingroup\small

\begin{longtable}[]{@{}
  >{\raggedright\arraybackslash}p{(\columnwidth - 4\tabcolsep) * \real{0.1296}}
  >{\raggedright\arraybackslash}p{(\columnwidth - 4\tabcolsep) * \real{0.4200}}
  >{\raggedright\arraybackslash}p{(\columnwidth - 4\tabcolsep) * \real{0.4504}}@{}}
\toprule\noalign{}
\begin{minipage}[b]{\linewidth}\raggedright
\textbf{Type}
\end{minipage} & \begin{minipage}[b]{\linewidth}\raggedright
\textbf{Key feature}
\end{minipage} & \begin{minipage}[b]{\linewidth}\raggedright
\textbf{Remember}
\end{minipage} \\
\midrule\noalign{}
\endhead
\bottomrule\noalign{}
\endlastfoot
Fundamental & Adds to knowledge/\allowbreak{}theory & Not immediately applicable \\
Applied & Solves immediate practical problem & Product-oriented \\
\textbf{Action} & Done by practitioner (teacher) to improve own practice; cyclic & Coined by Kurt Lewin; in education by Stephen Corey \\
Experimental & Manipulation of IV + control & Most scientific, cause-effect \\
Ex-post facto & IV already occurred, no manipulation & "After the fact" \\
Historical & Past events; primary \& secondary sources & External criticism (authenticity) \& internal criticism (credibility) \\
Descriptive & "What is" --- survey, correlational & No manipulation \\
Ethnography & Culture studied by immersion & Participant observation \\
Grounded theory & Theory emerges from data & Glaser \& Strauss \\
Phenomenology & Lived experiences & Husserl \\
\end{longtable}

\endgroup

\needspace{135pt}

\hypertarget{p1-02--positivism-vs-post-positivism}{%
\subsection{Positivism vs Post-positivism}\label{p1-02--positivism-vs-post-positivism}}

\begingroup\small

\begin{longtable}[]{@{}
  >{\raggedright\arraybackslash}p{(\columnwidth - 2\tabcolsep) * \real{0.2571}}
  >{\raggedright\arraybackslash}p{(\columnwidth - 2\tabcolsep) * \real{0.3302}}@{}}
\toprule\noalign{}
\begin{minipage}[b]{\linewidth}\raggedright
\textbf{Positivism (Auguste Comte)}
\end{minipage} & \begin{minipage}[b]{\linewidth}\raggedright
\textbf{Post-positivism / Interpretivism}
\end{minipage} \\
\midrule\noalign{}
\endhead
\bottomrule\noalign{}
\endlastfoot
Reality is objective, single & Reality is multiple, constructed \\
Quantitative, deductive & Qualitative (or critical realism) \\
Researcher detached & Researcher involved \\
Hypothesis testing, generalisation & Understanding meaning (\emph{Verstehen} -- Weber) \\
\end{longtable}

\endgroup

Paradigm (Thomas Kuhn) = ontology (nature of reality) + epistemology (nature of knowledge) + methodology.

\needdiagram{figures/p1-02-01}{55pt}

\hypertarget{p1-02--3-variables--hypothesis}{%
\section{3. Variables \& Hypothesis}\label{p1-02--3-variables--hypothesis}}

\diagramhere{figures/p1-02-01}

\begin{itemize}
\tightlist
\item
  \textbf{Hypothesis} = tentative, testable statement about relationship between variables.
\item
  \textbf{Null hypothesis (H₀)}: no relationship/difference. \textbf{Alternative (H₁)}: there is.
\item
  \textbf{Type I error (α)}: rejecting a TRUE null. \textbf{Type II error (β)}: accepting a FALSE null. Power = 1 − β.
\item
  Directional (one-tailed) vs non-directional (two-tailed).
\end{itemize}

\needdiagram{figures/p1-02-02}{55pt}

\hypertarget{p1-02--4-steps-of-research}{%
\section{4. Steps of Research}\label{p1-02--4-steps-of-research}}

\diagramhere{figures/p1-02-02}

\needdiagram{figures/p1-02-03}{55pt}

\hypertarget{p1-02--5-sampling}{%
\section{5. Sampling}\label{p1-02--5-sampling}}

\diagramhere{figures/p1-02-03}

\textbf{Tools of data collection}: questionnaire (closed/open), interview (structured, unstructured), observation
(participant, non-participant), schedule (filled by enumerator), rating scales (\textbf{Likert} -- 5-point agree/disagree;
\textbf{Thurstone} -- equal-appearing intervals; \textbf{Guttman} -- cumulative; \textbf{Semantic differential} -- Osgood; bipolar adjectives).

\textbf{Scales of measurement (Stevens)}: Nominal (labels) → Ordinal (rank) → Interval (no true zero, e.g. °C) → Ratio (true zero, e.g. weight).

\textbf{Parametric tests} (normal data): t-test, z-test, ANOVA (F), Pearson r.
\textbf{Non-parametric}: Chi-square, Mann-Whitney U, Wilcoxon, Kruskal-Wallis, Spearman rho, sign test.

\needspace{205pt}

\hypertarget{p1-02--6-thesis--article-writing}{%
\section{6. Thesis \& Article Writing}\label{p1-02--6-thesis--article-writing}}

Thesis format: Preliminary pages (title, declaration, certificate, acknowledgement, contents, list of tables) →
Main body (Introduction, Review of Literature, Methodology, Results, Discussion, Conclusion) → References → Appendices.

\textbf{IMRAD} structure of research article: Introduction, Methods, Results, And Discussion.

\begingroup\small

\begin{longtable}[]{@{}
  >{\raggedright\arraybackslash}p{(\columnwidth - 4\tabcolsep) * \real{0.3097}}
  >{\raggedright\arraybackslash}p{(\columnwidth - 4\tabcolsep) * \real{0.3214}}
  >{\raggedright\arraybackslash}p{(\columnwidth - 4\tabcolsep) * \real{0.2427}}@{}}
\toprule\noalign{}
\begin{minipage}[b]{\linewidth}\raggedright
\textbf{Style}
\end{minipage} & \begin{minipage}[b]{\linewidth}\raggedright
\textbf{Used in}
\end{minipage} & \begin{minipage}[b]{\linewidth}\raggedright
\textbf{Example in-text}
\end{minipage} \\
\midrule\noalign{}
\endhead
\bottomrule\noalign{}
\endlastfoot
\textbf{APA} (American Psychological Assoc.) & Social sciences, education & (Sharma, 2020) --- author-date \\
\textbf{MLA} (Modern Language Assoc.) & Humanities, literature & (Sharma 45) --- author-page \\
\textbf{Chicago} & History; notes-bibliography or author-date & Footnotes \\
\textbf{Harvard} & General author-date & (Sharma 2020) \\
\textbf{IEEE} & Engineering, CS & {[}1{]} numbered \\
\textbf{Vancouver} & Medicine & numbered \\
\end{longtable}

\endgroup

\begin{itemize}
\tightlist
\item
  \emph{Ibid.} = same source as previous note; \emph{Op. cit.} = work cited earlier (different page); \emph{Loc. cit.} = same place cited earlier; \emph{et al.} = and others.
\item
  \textbf{Bibliography} = all sources consulted; \textbf{References} = only those cited.
\item
  \textbf{Abstract}: \textasciitilde150--300 words, written last, placed first. \textbf{Annotated bibliography} includes short summaries.
\end{itemize}

\needspace{135pt}

\hypertarget{p1-02--7-ict-in-research}{%
\section{7. ICT in Research}\label{p1-02--7-ict-in-research}}

\begingroup\small

\begin{longtable}[]{@{}
  >{\raggedright\arraybackslash}p{(\columnwidth - 2\tabcolsep) * \real{0.1750}}
  >{\raggedright\arraybackslash}p{(\columnwidth - 2\tabcolsep) * \real{0.7934}}@{}}
\toprule\noalign{}
\begin{minipage}[b]{\linewidth}\raggedright
\textbf{Need}
\end{minipage} & \begin{minipage}[b]{\linewidth}\raggedright
\textbf{Tool}
\end{minipage} \\
\midrule\noalign{}
\endhead
\bottomrule\noalign{}
\endlastfoot
Literature search & Google Scholar, Scopus, Web of Science, PubMed, Shodhganga, INFLIBNET N-LIST, e-ShodhSindhu \\
Reference management & Mendeley, Zotero, EndNote \\
Statistical analysis & SPSS, R, SAS, STATA, Excel \\
Qualitative analysis & NVivo, ATLAS.ti, MAXQDA \\
Plagiarism detection & \textbf{Shodhshuddhi (DrillBit/\allowbreak{}Ouriginal-URKUND)}, Turnitin, iThenticate \\
Writing & LaTeX, MS Word \\
Unique researcher ID & ORCID, Scopus Author ID, ResearcherID \\
Metrics & Impact Factor (Clarivate JCR), h-index (Hirsch), i10-index (Google Scholar), CiteScore (Scopus) \\
\end{longtable}

\endgroup

\textbf{h-index}: a researcher has index h if h papers have ≥ h citations each.
\textbf{Impact factor (2024)} = citations in 2024 to items published in 2022--23 ÷ citable items published in 2022--23.

\needspace{100pt}

\hypertarget{p1-02--8-research-ethics}{%
\section{8. Research Ethics}\label{p1-02--8-research-ethics}}

\begin{itemize}
\tightlist
\item
  \textbf{Plagiarism}: UGC Regulations 2018 (Promotion of Academic Integrity and Prevention of Plagiarism):
\end{itemize}

\begingroup\small

\begin{longtable}[]{@{}
  >{\raggedright\arraybackslash}p{(\columnwidth - 4\tabcolsep) * \real{0.0723}}
  >{\raggedright\arraybackslash}p{(\columnwidth - 4\tabcolsep) * \real{0.1043}}
  >{\raggedright\arraybackslash}p{(\columnwidth - 4\tabcolsep) * \real{0.2942}}@{}}
\toprule\noalign{}
\begin{minipage}[b]{\linewidth}\raggedright
\textbf{Level}
\end{minipage} & \begin{minipage}[b]{\linewidth}\raggedright
\textbf{Similarity}
\end{minipage} & \begin{minipage}[b]{\linewidth}\raggedright
\textbf{Penalty (for students)}
\end{minipage} \\
\midrule\noalign{}
\endhead
\bottomrule\noalign{}
\endlastfoot
Level 0 & up to 10\% & Minor, no penalty \\
Level 1 & 10--40\% & Resubmit within 6 months \\
Level 2 & 40--60\% & Debarred from resubmission for 1 year \\
Level 3 & \textgreater{} 60\% & Registration cancelled \\
\end{longtable}

\endgroup

\begin{itemize}
\tightlist
\item
  Other ethics: informed consent, confidentiality, anonymity, no fabrication / falsification, honest authorship,
  avoid \textbf{predatory journals}; \textbf{UGC-CARE} list of quality journals (2019).
\item
  \textbf{Salami slicing}: splitting one study into many papers. \textbf{Duplicate publication} is unethical.
\item
  \textbf{COPE}: Committee on Publication Ethics.
\end{itemize}

\needspace{175pt}

\hypertarget{p1-02--9-deeper-dive--research-designs-reliability-validity--statistics}{%
\section{9. Deeper Dive --- Research Designs, Reliability, Validity \& Statistics}\label{p1-02--9-deeper-dive--research-designs-reliability-validity--statistics}}

\needspace{135pt}

\hypertarget{p1-02--91-experimental-designs-campbell--stanley-notation-r--random-assignment-o--observation-x--treatment}{%
\subsection{9.1 Experimental Designs (Campbell \& Stanley notation: R = random assignment, O = observation, X = treatment)}\label{p1-02--91-experimental-designs-campbell--stanley-notation-r--random-assignment-o--observation-x--treatment}}

\begingroup\small

\begin{longtable}[]{@{}
  >{\raggedright\arraybackslash}p{(\columnwidth - 4\tabcolsep) * \real{0.3738}}
  >{\raggedright\arraybackslash}p{(\columnwidth - 4\tabcolsep) * \real{0.2466}}
  >{\raggedright\arraybackslash}p{(\columnwidth - 4\tabcolsep) * \real{0.3204}}@{}}
\toprule\noalign{}
\begin{minipage}[b]{\linewidth}\raggedright
\textbf{Design}
\end{minipage} & \begin{minipage}[b]{\linewidth}\raggedright
\textbf{Notation}
\end{minipage} & \begin{minipage}[b]{\linewidth}\raggedright
\textbf{Notes}
\end{minipage} \\
\midrule\noalign{}
\endhead
\bottomrule\noalign{}
\endlastfoot
One-shot case study (pre-experimental) & X O & No comparison; weakest \\
One-group pre-test post-test & O X O & No control group \\
Static group comparison & X O / --- O & Non-random groups \\
\textbf{Pre-test post-test control group} (true) & R O X O / R O --- O & Classic true experiment \\
\textbf{Post-test only control group} & R X O / R --- O & Avoids testing effect \\
\textbf{Solomon four-group} & Combines both above (4 groups) & Controls pre-test sensitisation; strongest \\
Quasi-experimental: non-equivalent control group & O X O / O --- O & Intact groups (e.g. two existing classes) \\
Time-series & O O O X O O O & Repeated observations \\
Factorial design & 2 × 2, 2 × 3 \ldots{} & Two or more IVs and their interaction \\
\end{longtable}

\endgroup

\textbf{Threats to internal validity}: history, maturation, testing, instrumentation, statistical regression, selection bias, experimental mortality (attrition).
\textbf{External validity} = generalisability; threats: reactive effects (Hawthorne effect), interaction of selection and treatment.

\needdiagram{figures/p1-02-04}{55pt}

\hypertarget{p1-02--92-reliability--validity}{%
\subsection{9.2 Reliability \& Validity}\label{p1-02--92-reliability--validity}}

\diagramhere{figures/p1-02-04}

\begin{itemize}
\tightlist
\item
  Reliability is necessary but \textbf{not sufficient} for validity. A valid test must be reliable.
\item
  \textbf{Spearman--Brown}: full-test reliability r\_full = 2r\_half / (1 + r\_half).
\item
  Cronbach\textquotesingle s α ≥ 0.7 is usually acceptable.
\end{itemize}

\begin{asciiblock}{110pt}
\begin{Verbatim}[fontsize=\fontsize{8.8}{10.7}\selectfont]
 Target analogy
  Reliable, not valid      Valid & reliable        Neither
   ┌─────────────┐         ┌─────────────┐        ┌─────────────┐
   │   ○         │         │      ○      │        │ x        x  │
   │ xxx         │         │     xxx     │        │      ○      │
   │ xx          │         │     xx      │        │  x     x    │
   └─────────────┘         └─────────────┘        └─────────────┘
   (tight cluster off-centre) (cluster on centre)  (scattered)
\end{Verbatim}
\end{asciiblock}

\needspace{135pt}

\hypertarget{p1-02--93-statistics-for-researchers}{%
\subsection{9.3 Statistics for Researchers}\label{p1-02--93-statistics-for-researchers}}

\begingroup\small

\begin{longtable}[]{@{}
  >{\raggedright\arraybackslash}p{(\columnwidth - 2\tabcolsep) * \real{0.2160}}
  >{\raggedright\arraybackslash}p{(\columnwidth - 2\tabcolsep) * \real{0.6406}}@{}}
\toprule\noalign{}
\begin{minipage}[b]{\linewidth}\raggedright
\textbf{Concept}
\end{minipage} & \begin{minipage}[b]{\linewidth}\raggedright
\textbf{Key point}
\end{minipage} \\
\midrule\noalign{}
\endhead
\bottomrule\noalign{}
\endlastfoot
Correlation r (Pearson) & −1 ≤ r ≤ +1; r = 0 → no linear relation; correlation ≠ causation \\
Spearman rank ρ & ρ = 1 − 6Σd² / (n(n² − 1)) for ordinal data \\
Coefficient of determination & r² = proportion of variance explained \\
Regression & Predicts Y from X: Y = a + bX \\
Level of significance & α = 0.05 or 0.01 (probability of Type I error) \\
Degrees of freedom & t-test (one sample) n − 1; chi-square (r − 1)(c − 1) \\
Normal curve & Mean = median = mode; 68.26\% within ±1σ, 95.44\% within ±2σ, 99.74\% within ±3σ \\
Skewness & Positive: mean \textgreater{} median \textgreater{} mode; Negative: mean \textless{} median \textless{} mode \\
Kurtosis & Leptokurtic (peaked), Mesokurtic (normal), Platykurtic (flat) \\
\end{longtable}

\endgroup

\textbf{Worked (Spearman ρ)}: n = 5, Σd² = 4 → ρ = 1 − (6 × 4)/(5 × 24) = 1 − 0.2 = \textbf{0.8}.

\needspace{135pt}

\hypertarget{p1-02--94-research-proposal-synopsis-structure}{%
\subsection{9.4 Research Proposal (Synopsis) Structure}\label{p1-02--94-research-proposal-synopsis-structure}}

Title → Introduction \& background → Review of literature \& research gap → Statement of the problem → Objectives → Hypotheses → Methodology (design, population, sample, tools, analysis) → Delimitations → Significance → Chapterisation → Time-frame → References.

\begin{itemize}
\tightlist
\item
  \textbf{Delimitations} = boundaries set by the researcher; \textbf{Limitations} = constraints beyond the researcher\textquotesingle s control.
\item
  \textbf{Operational definition} = defining a variable in terms of how it is measured.
\end{itemize}

\needspace{135pt}

\hypertarget{p1-02--previous-year-questions-pyq-pattern}{%
\section{Previous Year Questions (PYQ pattern)}\label{p1-02--previous-year-questions-pyq-pattern}}

\textbf{Types of research}

\begin{enumerate}
\def\labelenumi{\arabic{enumi}.}
\tightlist
\item
  A teacher conducts research to solve a problem of absenteeism in her class. It is: \textbf{Action research}
\item
  Research aimed at developing theory is: \textbf{Fundamental}
\item
  Study of the effect of an event that has already happened without manipulation: \textbf{Ex-post facto}
\item
  "Who coined action research?" --- \textbf{Kurt Lewin} (in education popularised by Stephen Corey)
\item
  In historical research, testing authenticity of a source is: \textbf{External criticism}
\end{enumerate}

\textbf{Positivism}

\begin{enumerate}
\def\labelenumi{\arabic{enumi}.}
\setcounter{enumi}{5}
\tightlist
\item
  Positivism was proposed by: \textbf{Auguste Comte}
\item
  Which is associated with qualitative research? (a) generalisation (b) hypothesis testing (c) thick description (d) large samples --- \textbf{Ans: (c)}
\end{enumerate}

\textbf{Hypothesis \& errors}

\begin{enumerate}
\def\labelenumi{\arabic{enumi}.}
\setcounter{enumi}{7}
\tightlist
\item
  Rejecting a true null hypothesis is: \textbf{Type I error}
\item
  Hypothesis that states "no difference" is: \textbf{Null hypothesis}
\item
  The variable manipulated by the researcher: \textbf{Independent variable}
\end{enumerate}

\textbf{Sampling}

\begin{enumerate}
\def\labelenumi{\arabic{enumi}.}
\setcounter{enumi}{10}
\tightlist
\item
  Studying drug users via referral chains: \textbf{Snowball sampling}
\item
  Population divided into homogeneous subgroups then randomly sampled: \textbf{Stratified random}
\item
  Which is a non-probability method? (a) cluster (b) systematic (c) quota (d) stratified --- \textbf{Ans: (c)}
\end{enumerate}

\textbf{Scales \& tools}

\begin{enumerate}
\def\labelenumi{\arabic{enumi}.}
\setcounter{enumi}{13}
\tightlist
\item
  Temperature in °C is measured on which scale? \textbf{Interval}
\item
  A 5-point agree--disagree scale is: \textbf{Likert scale}
\item
  Chi-square test is: \textbf{Non-parametric}
\end{enumerate}

\textbf{Writing \& referencing}

\begin{enumerate}
\def\labelenumi{\arabic{enumi}.}
\setcounter{enumi}{16}
\tightlist
\item
  "Ibid." is used when: \textbf{citing the same source as the immediately preceding citation}
\item
  APA style follows: \textbf{author--date system}
\item
  The abstract of a thesis is written: \textbf{at the end, placed at the beginning}
\item
  Correct sequence of research steps: \textbf{Problem → Literature → Hypothesis → Design → Data → Analysis → Report}
\end{enumerate}

\textbf{ICT \& Ethics}

\begin{enumerate}
\def\labelenumi{\arabic{enumi}.}
\setcounter{enumi}{20}
\tightlist
\item
  UGC\textquotesingle s plagiarism software made available to universities: \textbf{Shodhshuddhi (URKUND/Ouriginal, later DrillBit)}
\item
  As per UGC 2018 regulations, similarity up to 10\% is: \textbf{Level 0 (no penalty)}
\item
  h-index was proposed by: \textbf{J.E. Hirsch (2005)}
\item
  Which identifier uniquely identifies a researcher? \textbf{ORCID}
\item
  Reference management software: \textbf{Mendeley / Zotero}
\end{enumerate}

\textbf{More practice questions}

\begin{enumerate}
\def\labelenumi{\arabic{enumi}.}
\setcounter{enumi}{25}
\tightlist
\item
  The design that controls the effect of pre-testing using four groups: \textbf{Solomon four-group design}
\item
  Using two intact classes without random assignment is a: \textbf{quasi-experimental design}
\item
  Improvement in performance because subjects know they are being observed: \textbf{Hawthorne effect}
\item
  Consistency of scores over time is measured by: \textbf{test-retest reliability}
\item
  Cronbach\textquotesingle s alpha measures: \textbf{internal consistency}
\item
  Split-half reliability is corrected by: \textbf{Spearman--Brown formula}
\item
  A test predicting future job performance has: \textbf{predictive validity}
\item
  In a normal distribution, the percentage of cases within ±1σ: \textbf{68.26\%}
\item
  When mean \textgreater{} median \textgreater{} mode, distribution is: \textbf{positively skewed}
\item
  Degrees of freedom for a 3 × 4 contingency table chi-square: \textbf{6}
\item
  Boundaries set by the researcher on the scope of study: \textbf{delimitations}
\item
  r = −0.9 indicates: \textbf{strong negative correlation}
\item
  Which is NOT a threat to internal validity? (a) history (b) maturation (c) generalisability (d) attrition --- \textbf{Ans: (c)}
\item
  Defining "academic achievement" as "marks obtained in the annual exam" is an: \textbf{operational definition}
\end{enumerate}

\begin{revbox}

\begin{itemize}
\tightlist
\item
  Action research → practitioner, local, immediate · Ex-post facto → no manipulation
\item
  Type I = reject true H₀ (α) · Type II = accept false H₀ (β)
\item
  NOIR = Nominal, Ordinal, Interval, Ratio
\item
  Plagiarism levels: 0 (≤10), 1 (10--40), 2 (40--60), 3 (\textgreater60)
\end{itemize}

\end{revbox}
